\documentclass[12pt,reqno]{amsart}
\usepackage{math-review}
% Real Banach spaces; all time norms on a fixed finite interval.
\newcommand{\trans}{\tau}
\newcommand{\cone}{\mathcal M_{+}}
\newcommand{\gauge}[1]{[#1]_{\cone}}
\newcommand{\R}{\mathbb R}

\title{Time compactness: hypotheses, failures and exact overlap}
\date{}
\begin{document}
\maketitle

% REGION result-map BEGIN
Compactness can fail because the wrong unknown is tracked, the time exponent is too weak, or an algebraic operation leaves the controlled set. Four local units distinguish actual counterexamples from limits of a proof or a change of variables.
% REGION result-map END

% REGION conventions BEGIN
All spaces and the finite interval $(0,T)$ are fixed. Relative compactness means that every sequence has a norm-convergent subsequence. Write
\[
\tau_h f(t)=f(t+h),\qquad 0<t<T-h.
\]
Every translation limit below is uniform over the stated family. At $p=\infty$, an $L^\infty$ class needs a continuous representative before it can be treated in $C([0,T];E)$.

We assume Bochner integration, H\"older and Minkowski inequalities, completeness, and the fundamental theorem for $W^{1,1}$ trajectories. The compactness criterion is imported precisely; the nonlinear constructions are developed here.
% REGION conventions END

\section{Which unknown is compact?}

% REGION N-maitre BEGIN
\begin{theorem}[Maitre]\label{thm:maitre}
Let $T>0$, $1\le p\le\infty$, and let $E_1,E_2$ be Banach spaces.
Let $K:E_1\to E_2$ map every bounded subset of $E_1$ into a relatively compact subset of $E_2$.
Let $U\subset L^1(0,T;E_1)$ be bounded and suppose that each $v(t)=K(u(t))$, $u\in U$, is strongly measurable.
Write $V=\{K(u):u\in U\}$ and assume $V\subset L^p(0,T;E_2)$, that for some $1<r\le\infty$
\[
\sup_{v\in V}\|v\|_{L^r(0,T;E_2)}<\infty,
\qquad
\lim_{h\downarrow0}\sup_{v\in V}\|v(\cdot+h)-v\|_{L^p(0,T-h;E_2)}=0.
\]
Then $V$ is relatively compact in $L^p(0,T;E_2)$ for $p<\infty$.
For $p=\infty$, the continuous representatives form a relatively compact set in $C([0,T];E_2)$.
No injectivity, linearity or monotonicity of $K$ is required. The conclusion concerns $K(u)$, not $u$.
\end{theorem}
% REGION N-maitre END

% REGION map-orientation BEGIN
The conclusion of \cite[Theorem 2.1, p.~1726]{S-maitre} is for $K(u)$. No inverse, linearity or monotonicity is assumed. The following example shows why replacing $K(u)$ by $u$ changes the claim.
% REGION map-orientation END

% REGION image-boundary BEGIN
Take $E_1=E_2=\mathbb R$, $K(x)=x^2$, and
\[
u_n(t)=(-1)^{\lfloor2^nt\rfloor},\qquad 0<t<1.
\]
For $m>n$, each dyadic cell for $u_n$ contains an even number of alternating cells for $u_m$. The signs differ on exactly half the interval, so
\[
\|u_n-u_m\|_{L^1}=2\cdot\tfrac12=1,
\qquad K(u_n)=1.
\]
The inputs have no $L^1$-Cauchy subsequence, although their images are a singleton with zero translations. This tests the exact unknown in \cite[Theorem 2.1 and Remark 2.2(3), p.~1726]{S-maitre}.

It also prevents using a formal $K^{-1}$ for every compact map. The example contradicts input recovery; it fully agrees with compactness of the images.
% REGION image-boundary END
\subsection{Where the integrability premise enters}

% REGION truncation-tail BEGIN
Let $C_U=\sup_{u\in U}\|u\|_{L^1(E_1)}$ and $C_V=\sup_{v\in V}\|v\|_{L^r(E_2)}$. For $M>0$, define
\[
G_M(u)=\{t:\|u(t)\|_{E_1}>M\},\qquad
u_M=u\mathbf1_{G_M(u)^c},\qquad v_M=K(u_M).
\]
Markov bounds the exceptional set. Since $v_M$ equals either $v$ or the fixed value $K(0)$, it is measurable and
\[
|G_M(u)|\le C_U/M,\qquad v-v_M=\mathbf1_{G_M(u)}(v-K(0)).
\]
H\"older therefore gives the uniform tail estimate
\begin{equation}\label{eq:tail}
\|v-v_M\|_{L^1(E_2)}
\le (C_U/M)^{1-1/r}C_V+(C_U/M)\|K(0)\|_{E_2}.
\end{equation}
The term $K(0)$ cannot be dropped for an arbitrary compact map. The right side tends to zero for $r>1$, with $1/r=0$ at $r=\infty$.

Only an $L^1$ error is needed: it controls the error of every interval integral. We do not infer a small $L^p$ tail when $p>r$.
% REGION truncation-tail END

% REGION r-limit BEGIN
At $r=1$, the power $1-1/r$ in \eqref{eq:tail} is zero. This bound no longer proves uniform smallness of the discarded output. It does not prove that all nonlinear compactness conclusions fail at $r=1$.

The account is restricted to these three source interfaces. Later nonlinear translation refinements are outside its scope; no optimality or strongest-known claim is made for the auxiliary $r>1$ premise.
% REGION r-limit END
\subsection{Why the valid conclusion still follows}

% REGION N-simon-criterion BEGIN
\begin{lemma}[Simon]\label{lem:criterion}
Let $E$ be a Banach space, $T>0$, $1\le p\le\infty$, and $F\subset L^p(0,T;E)$.
Assume that for every $0<a<b<T$ the set
\[
\left\{\int_a^b f(t)\,dt:f\in F\right\}
\]
is relatively compact in $E$, and that
\[
\lim_{h\downarrow0}\sup_{f\in F}\|f(\cdot+h)-f\|_{L^p(0,T-h;E)}=0.
\]
Then $F$ is relatively compact in $L^p(0,T;E)$ if $p<\infty$.
If $p=\infty$, its elements have continuous representatives on $[0,T]$ and these representatives form a relatively compact subset of $C([0,T];E)$ with the supremum norm.
\end{lemma}
% REGION N-simon-criterion END

% REGION criterion-use BEGIN
This is the sufficient direction of \cite[Theorem 1, pp.~71--73]{S-simon}. Its proof uses right time averages and then reverses time to cover the other endpoint.

It reduces each route to two joint tasks: compact interval integrals and small translations in the same ambient norm. At infinity it also supplies continuous representatives; it does not characterize all compact subsets of $L^\infty$.
% REGION criterion-use END

% REGION compact-hull BEGIN
For fixed $M$, all values of $v_M$ lie in the compact set
\[
Q_M=\overline{K(\{x:\|x\|_{E_1}\le M\})}.
\]
To average such values, take a finite $\delta$-net $q_1,\ldots,q_J$ of a compact set $Q\subset E$. Replacing each point of a convex combination by a nearby net point shows
\[
\operatorname{co}(Q)\subset
\operatorname{co}\{q_1,\ldots,q_J\}+\delta B_E.
\]
The finite-dimensional simplex is compact. Thus the hull is totally bounded, and its closure is compact by completeness.

For measurable $w:I\to Q$, use the first nearest net point to break ties. The resulting cells are measurable. Their relative measures are nonnegative weights summing to one, so
\[
|I|^{-1}\int_Iw(t)\,dt\in\overline{\operatorname{co}}(Q).
\]
Indeed, the simple-function average lies in the finite simplex and its error is at most $\delta$. Letting $\delta$ decrease gives the displayed inclusion.

Hence $\{\int_Iv_M\}$ is relatively compact. Choose $M$ to control \eqref{eq:tail}, then a finite net for this fixed truncated set. The original interval integrals are totally bounded in $E_2$.

Lemma~\ref{lem:criterion} now uses the assumed $L^p(E_2)$ translations to conclude Theorem~\ref{thm:maitre}. Maitre uses shifted Riemann sums \cite[Lemma 2.3, pp.~1726--1730]{S-maitre}; the compact-hull proof here is a local reconstruction.
% REGION compact-hull END

\section{How weak may the time control be?}

% REGION N-simon BEGIN
\begin{theorem}[Simon]\label{thm:simon}
Let $T>0$, $1\le p\le\infty$, and let $X,B,Y$ be Banach spaces with continuous linear injections
\[
X\Subset B\hookrightarrow Y,
\]
where the first injection is compact. Let $F\subset L^p(0,T;X)$ satisfy
\[
\sup_{f\in F}\|f\|_{L^p(0,T;X)}<\infty,
\qquad
\lim_{h\downarrow0}\sup_{f\in F}\|f(\cdot+h)-f\|_{L^p(0,T-h;Y)}=0.
\]
Then $F$ is relatively compact in $L^p(0,T;B)$ for $p<\infty$.
For $p=\infty$, the $B$-valued continuous representatives of the elements of $F$ form a relatively compact set in $C([0,T];B)$.
\end{theorem}
% REGION N-simon END

% REGION linear-orientation BEGIN
The weak space $Y$ in \cite[Theorem 5, pp.~84--85]{S-simon} is allowed because the compact spatial estimate transfers norms. Translation continuity of each individual trajectory would not supply the uniform modulus in the statement.
% REGION linear-orientation END

% REGION linear-proof BEGIN
Put $M_X=\sup_{f\in F}\|f\|_{L^p(X)}$. For each interior interval $I$, H\"older gives
\[
\left\|\int_I f(t)\,dt\right\|_X\le |I|^{1-1/p}M_X.
\]
These integrals are relatively compact in $B$ and hence in $Y$. Lemma~\ref{lem:criterion}, applied in $Y$, gives an $L^p(Y)$-Cauchy subsequence.

The remaining operation changes norms. Compactness and injectivity yield the Ehrling estimate \cite[Lemma 8, pp.~84--85]{S-simon}:
\begin{equation}\label{eq:linear-ehrling}
\|z\|_B\le\eta\|z\|_X+C_\eta\|z\|_Y,\qquad z\in X,\quad\eta>0.
\end{equation}
Indeed, failure at $\eta_0$ gives a sequence normalized to unit $X$ norm, with $B$ norm above $\eta_0$ and $Y$ norm tending to zero. Compactness gives a $B$ limit; injectivity makes it zero, a contradiction.

Apply \eqref{eq:linear-ehrling} to $f_n(t)-f_m(t)$ and use Minkowski:
\begin{equation}\label{eq:linear-upgrade}
\|f_n-f_m\|_{L^p(B)}
\le2\eta M_X+C_\eta\|f_n-f_m\|_{L^p(Y)}.
\end{equation}
For $M_X>0$, choose $\eta<\varepsilon/(4M_X)$ first. Then choose $n,m$ to make the second term smaller than $\varepsilon/2$. If $M_X=0$, the family vanishes.

Completeness gives the $L^p(B)$ limit. At infinity, on the full-measure dense set where $\|f(t)\|_X\le M_X$, the same estimate gives uniform $B$ continuity from the continuous $Y$ representative.

First choose $\eta$, then $|t-s|$; extend to $[0,T]$ by completeness. The injection identifies the extension with the $Y$ representative. Thus the supremum conclusion in Theorem~\ref{thm:simon} concerns actual continuous functions.
% REGION linear-proof END
\subsection{The derivative exponent and its actual obstruction}

% REGION derivative-modulus BEGIN
For $v\in W^{1,1}(0,T;E)$, choose a time with norm no larger than the time average plus an arbitrary error. Integrate $v'$ from that time and let the error vanish:
\begin{equation}\label{eq:w11}
\|v\|_{L^\infty(E)}\le T^{-1}\|v\|_{L^1(E)}+\|v'\|_{L^1(E)}.
\end{equation}
If the resulting amplitude bound is $K_0$ and $\|v'\|_{L^1(E)}\le D$, Fubini gives
\begin{equation}\label{eq:l1shift}
\|\tau_hv-v\|_{L^1(E)}
\le\int_0^{T-h}\int_t^{t+h}\|v'(s)\|_E\,ds\,dt\le hD.
\end{equation}
The pointwise difference is at most $2K_0$. Hence, for every finite $p\ge1$,
\begin{equation}\label{eq:pshift}
\|\tau_hv-v\|_{L^p(E)}^p
\le(2K_0)^{p-1}\|\tau_hv-v\|_{L^1(E)}
\le(2K_0)^{p-1}hD.
\end{equation}
For $q>1$, H\"older on $(t,t+h)$ instead supplies a supremum modulus:
\begin{equation}\label{eq:supshift}
\|v(t+h)-v(t)\|_E\le h^{1-1/q}\|v'\|_{L^q(E)}.
\end{equation}
Take $1/q=0$ for $q=\infty$. At $q=1$, this last power vanishes. The finite-$p$ estimate has order $h^{1/p}$, not $h$; compare \cite[Lemma 4, pp.~77--78]{S-simon} and \cite[Corollary 2.4, p.~1730]{S-maitre}.
% REGION derivative-modulus END

% REGION endpoint-boundary BEGIN
On $(0,1)$, let $n\ge3$ and set
\[
v_n(t)=\begin{cases}
0,&t\le\tfrac12,\\
n(t-\tfrac12),&\tfrac12<t<\tfrac12+\tfrac1n,\\
1,&t\ge\tfrac12+\tfrac1n.
\end{cases}
\]
These are continuous trajectories in the cone $[0,\infty)\subset\mathbb R$, with gauge $[z]=|z|$. Their spatial bounds are harmless, and
\[
\|v_n\|_\infty=1,\qquad \|v_n'\|_1=n\cdot n^{-1}=1.
\]
Their pointwise limit is the discontinuous step $w=\mathbf1_{(1/2,1]}$. No subsequence converges uniformly, since a uniform limit of continuous functions is continuous.

For continuous differences, the supremum and essential supremum agree. An $L^\infty$-Cauchy subsequence would therefore be uniformly Cauchy and have a continuous uniform limit, the same contradiction.

Yet every finite exponent gives convergence:
\begin{equation}\label{eq:ramp}
\|v_n-w\|_p^p=\int_0^{1/n}(1-ns)^p\,ds
=\frac1{n(p+1)}\longrightarrow0.
\end{equation}
Thus $(p,q)=(\infty,1)$ is a genuine obstruction already in one dimension; compare \cite[Proposition 3, p.~94]{S-simon} and \cite[Remark 2, p.~4]{S-barrett-suli}.
% REGION endpoint-boundary END

\section{Why does a cone require two points?}

% REGION N-cone BEGIN
\begin{theorem}[Barrett--S\"uli]\label{thm:cone}
Let $T>0$, let $A_0,A_1$ be real Banach spaces with a continuous linear injection $A_0\hookrightarrow A_1$, and let $\mathcal M_+\subset A_0$ be closed under multiplication by nonnegative scalars.
Suppose its gauge $[\cdot]:\mathcal M_+\to[0,\infty)$ satisfies
\[
[z]=0\Longleftrightarrow z=0,\qquad [cz]=c[z]\quad(c\ge0),
\qquad \|z\|_{A_0}\le C_0[z],
\]
and every gauge-bounded sequence in $\mathcal M_+$ has a subsequence converging in $A_0$.
Fix $1\le p,q\le\infty$ with either $q>1$, or $q=1$ and $p<\infty$.
Let $F$ be a family of strongly measurable $A_0$-valued functions $v$ on $(0,T)$ with $v(t)\in\mathcal M_+$ almost everywhere and measurable $t\mapsto[v(t)]$.
Assume the distributional derivatives $v'$ are $A_1$-valued and
\[
\sup_{v\in F}\left(\|[v(\cdot)]\|_{L^p(0,T)}+
\|v'\|_{L^q(0,T;A_1)}\right)<\infty.
\]
Then $F$ is relatively compact in $L^p(0,T;A_0)$.
The assertion includes $p=\infty$ only when $q>1$; it does not assert that limits belong to a nonclosed cone.
\end{theorem}
% REGION N-cone END

% REGION cone-orientation BEGIN
The cone theorem \cite[Theorem 1 and Remark 2, pp.~3--4]{S-barrett-suli} excludes exactly the pair tested above. A separate issue is algebraic: $[u-v]$ may be undefined, even when both gauges $[u]$ and $[v]$ are finite.
% REGION cone-orientation END

% REGION N-cone-ehrling BEGIN
\begin{lemma}[Barrett--S\"uli]\label{lem:ehrling}
Let $A_0\hookrightarrow A_1$ be a continuous linear injection of real Banach spaces.
Let $\mathcal M_+\subset A_0$ be a seminormed nonnegative cone with gauge $[\cdot]$, satisfying $\|z\|_{A_0}\le C[z]$ and compactness of gauge-bounded sequences in $A_0$.
Then for every $\eta>0$ there is $C_\eta<\infty$ such that
\[
\|u-v\|_{A_0}\le\eta([u]+[v])+C_\eta\|u-v\|_{A_1}
\quad\text{for all }u,v\in\mathcal M_+.
\]
Here the cone is closed under nonnegative scalar multiplication and its gauge is positive definite and positively degree-one homogeneous; no addition or subtraction closure is assumed.
\end{lemma}
% REGION N-cone-ehrling END

% REGION two-point-proof BEGIN
The proof in \cite[Lemma 1, pp.~4--5]{S-barrett-suli} explains why the sum of two gauges appears. If the estimate fails at $\eta_0>0$, choose $u_n,v_n$ with
\[
\|u_n-v_n\|_{A_0}>
\eta_0([u_n]+[v_n])+n\|u_n-v_n\|_{A_1}.
\]
The denominator $a_n=[u_n]+[v_n]$ is positive. Set $\widetilde u_n=u_n/a_n$ and $\widetilde v_n=v_n/a_n$. Each normalized gauge is at most one.

A common subsequence converges in $A_0$. Its difference is bounded there, while
\[
\|\widetilde u_n-\widetilde v_n\|_{A_0}>\eta_0,
\qquad
\|\widetilde u_n-\widetilde v_n\|_{A_1}
<n^{-1}\|\widetilde u_n-\widetilde v_n\|_{A_0}\longrightarrow0.
\]
Continuity and injectivity force the two $A_0$ limits to agree, contradicting the separation. Scaling each point stays inside the cone; taking a gauge of their difference would be illegitimate.
% REGION two-point-proof END
\subsection{Use only the operations permitted by the hypotheses}

% REGION cone-proof BEGIN
Write $G=\sup_{v\in F}\|[v]\|_{L^p}$, $D=\sup_{v\in F}\|v'\|_{L^q(A_1)}$, and $C_1=\|A_0\hookrightarrow A_1\|$. The gauge controls the weak-space time average:
\[
\|v\|_{L^1(A_1)}\le C_1C_0T^{1-1/p}G.
\]
Equation~\eqref{eq:w11} then gives
\begin{equation}\label{eq:coneamp}
\|v\|_{L^\infty(A_1)}\le K_1:=C_1C_0T^{-1/p}G+T^{1-1/q}D.
\end{equation}
This supplies continuous $A_1$ representatives. It does not require cone membership at every time or closedness of the cone.

Cut off by the gauge: $E_R(v)=\{t:[v(t)]>R\}$ and $v_R=v\mathbf1_{E_R(v)^c}$. Markov and \eqref{eq:coneamp} give
\begin{equation}\label{eq:conetail}
|E_R(v)|\le\frac{T^{1-1/p}G}{R},\qquad
\|v-v_R\|_{L^1(A_1)}\le\frac{K_1T^{1-1/p}G}{R}.
\end{equation}
The truncated values have a compact closure in $A_0$ and hence $A_1$. The compact-hull argument already proved gives compact interval integrals in $A_1$; the uniform tail passes this to $v$.

If $q=1$ and $p<\infty$, use \eqref{eq:pshift} in $A_1$. If $q>1$, use \eqref{eq:supshift}. Lemma~\ref{lem:criterion} gives a subsequence Cauchy in $L^p(A_1)$ in every allowed case.

Finally, Lemma~\ref{lem:ehrling} and Minkowski yield
\begin{equation}\label{eq:coneupgrade}
\|v_n-v_m\|_{L^p(A_0)}
\le2\eta G+C_\eta\|v_n-v_m\|_{L^p(A_1)}.
\end{equation}
For $G>0$, choose $\eta<\varepsilon/(4G)$ before the indices. Completeness gives strong $L^p(A_0)$ convergence. The case $G=0$ is immediate because every trajectory vanishes almost everywhere.

At $p=\infty$, the full-measure dense set with $[v(t)]\le G$ also gives continuous $A_0$ representatives. Apply two-point Ehrling to times $s,t$: choose $\eta$, then $|t-s|$, and extend by completeness.

The upgrade agrees with \cite[equation (2.22), pp.~12--13]{S-barrett-suli}. The weak-space argument here avoids the unproved time-dependent-index step following its (2.8), p.~6. This scoped concern is not a claim that the theorem is false.
% REGION cone-proof END

\section{When may one formulation replace another?}

% REGION N-homogeneous-transfer BEGIN
\begin{proposition}\label{prop:overlap}
Let $E_1,E_2$ be real Banach spaces, $T>0$, and $\alpha\ge1$.
Suppose $K:E_1\to E_2$ is injective, maps bounded sets into relatively compact sets, and satisfies $K(cu)=c^\alpha K(u)$ for $c\ge0$.
Set $\mathcal M_+=K(E_1)$ and $[v]=\|K^{-1}v\|_{E_1}^{\alpha}$.
Let $1\le p,q\le\infty$, with $q>1$ or with $q=1,p<\infty$.
Suppose a family $F$ of strongly measurable $\mathcal M_+$-valued trajectories has strongly measurable inverse lifts $u=K^{-1}v$, measurable gauges, and
\[
\sup_{v\in F}\bigl(\|[v]\|_{L^p(0,T)}+\|v'\|_{L^q(0,T;E_2)}\bigr)<\infty,
\]
where derivatives are $E_2$-valued distributions.
Then the hypotheses of the compact nonlinear-map theorem stated for $K$, $U=K^{-1}F$, and $V=F$ hold, using any fixed $r>1$ and the target exponent $p$.
Consequently $F$ is relatively compact in $L^p(0,T;E_2)$, and in $C([0,T];E_2)$ when $q>1$.
This reduction concerns the special cone $K(E_1)$ with the same output and derivative space $E_2$; it is not an equivalence for arbitrary cones or arbitrary compact maps.
\end{proposition}
% REGION N-homogeneous-transfer END

% REGION overlap-proof BEGIN
The special reduction in \cite[Section 4, pp.~15--16]{S-barrett-suli} can be checked directly. Injectivity and homogeneity give $K(0)=0$ and
\[
K^{-1}(cv)=c^{1/\alpha}K^{-1}v,\qquad [cv]=c[v].
\]
The compact image of the unit ball is bounded. Scaling therefore gives
\[
\|K(x)\|_{E_2}\le C\|x\|_{E_1}^{\alpha}=C[K(x)].
\]
Gauge-bounded sets come from bounded input sets and are relatively compact in $E_2$. This verifies the cone axioms and compact embedding.

Let $G=\sup\|[v]\|_{L^p}$. Since $\alpha p\ge1$, H\"older gives
\begin{equation}\label{eq:lift}
\|K^{-1}v\|_{L^1(E_1)}\le T^{1-1/(\alpha p)}G^{1/\alpha},
\end{equation}
with factor $T$ at $p=\infty$. The assumed measurable inverse lifts thus form a bounded $L^1(E_1)$ family.

The gauge and derivative bounds give a uniform $W^{1,1}(E_2)$ bound. Equation~\eqref{eq:w11} supplies every fixed $L^r(E_2)$ bound, $r>1$; \eqref{eq:pshift} or \eqref{eq:supshift} supplies the required translations.

Theorem~\ref{thm:maitre} applies, using $p=\infty$ when $q>1$ to obtain the $C$ conclusion. The derivative must be measured in $E_2$, the output space. A derivative in a larger weak space cannot simply be substituted.

Nor does an arbitrary cone come with this parametrization. Injectivity is needed for this inverse gauge; it is not an assumption of Maitre's theorem. These are limits of the reduction, not proofs that every alternative construction fails.
% REGION overlap-proof END

% REGION boundaries-synthesis BEGIN
The three restrictions have different logical force. Dyadic signs disprove recovery of input compactness. The ramp disproves adding the $(\infty,1)$ endpoint. By contrast, the failed $r=1$ tail estimate only identifies an unproved step.

Likewise, missing injectivity or a mismatched derivative space blocks the displayed inverse-gauge reduction. It does not assert that no other proof or parametrization could exist. A comparison must preserve the unknown, norm, exponents and quantifiers.
% REGION boundaries-synthesis END

% REGION source-scope BEGIN
The controlling texts are Simon's journal pagination and DOI, Maitre's published 2003 article, and Barrett--S\"uli's arXiv v2 of 16 April 2011. The latter's different internal 2018 date is not used to infer a new version.

The argument imports Simon's criterion and reconstructs the selected mechanisms. It claims no original research result, external referee certification or real-reader study. The historical addition of the cone source is preserved separately from these new readings.
% REGION source-scope END

\bibliographystyle{plain}
\bibliography{../generated/references}
\end{document}
